\section{Alternative geometric and entropy arguments}
\label{a:section}

We give two other descriptions of supported face codimension: a laminar
family of tight cuts, and a quotient-space formulation of contraction.
Both recover the bound used in the main proof. We then retain the direct
coefficient-three and coefficient-four covariance comparisons and their
counting consequences. These comparisons use weaker dimension estimates;
they do not require the finite-order entropy theorem.

Throughout, \(G=(V,E)\) is a finite loopless labelled multigraph on
\(n=2m\ge2\) vertices with a nonempty matching polytope, denoted by \(P=P(G)\). We use
\(F_x\), \(T(x)\), and \(d(x)=\dim T(x)\) from
Lemma~\ref{geom:lem-face}. Its direction-space description and
Theorem~\ref{geom:thm-polytope} are the shared geometric preliminaries.
Removing the zero coordinates of a mean preserves the dimension of its
minimal face: coordinate restriction and zero extension are inverse
affine maps on that face. We use this identification whenever we work
only on supported edges.

\subsection{Uncrossing and a laminar rank bound}

This alternative proof uses the laminar uncrossing method for tight odd cuts,
as in \cite[full version, Lemma~2.2 and Appendix~A]{SvenssonTarnawski2017};
compare also the tight-cut rank estimate in
\cite[Lemma~6.5]{AbdiCornuejolsDadushDalirrooyfard2026}. We give the spanning argument and the
tree count explicitly, retaining the dependence between two
complementary children of the root.

Recall that \(\chi_S=\mathbf1_{\delta(S)}\) and that
\(\mathcal S_x\) is the family of odd sets whose cuts are tight at
\(x\). Equation~\eqref{geom:eq-rank} identifies the supported face
codimension with the rank of these cut vectors after zero coordinates
are removed. We first choose a
laminar spanning family, then bound its rank by a tree count.

A family of sets is \emph{laminar} if every two of its members are
disjoint or one contains the other. Two sets \emph{cross} if they
intersect and neither contains the other.

\begin{lemma}[Laminar spanning family]
\label{geom:lem-laminar}
Suppose \(x\in P(G)\) and \(x_e>0\) for every edge. There is a
laminar family \(\mathcal L\subseteq\mathcal S_x\) whose cut
vectors are linearly independent and span all the cut vectors
\(\chi_S\), \(S\in\mathcal S_x\).
\end{lemma}

\begin{proof}
First let \(A,B\in\mathcal S_x\) cross, and put
\[
 I=A\cap B,\qquad D=A\setminus B,\qquad
 D'=B\setminus A,\qquad O=V\setminus(A\cup B).
\]
The first three regions are nonempty; \(O\) may be empty. If
\(|I|\) is odd, then both \(I\) and \(A\cup B\) are odd, and
the coordinatewise cut identity is
\begin{equation}
\label{geom:eq-uncross-intersection}
 \chi_A+\chi_B=
 \chi_I+\chi_{A\cup B}+2\mathbf1_{E(D,D')}.
\end{equation}
Here \(E(D,D')\) denotes the edges joining those two regions.
Taking the dot product with \(x\), the left side equals two, the
first two terms on the right are at least one each, and the last
term is nonnegative. Thus both new odd sets are tight and
\(x(E(D,D'))=0\). Since every coordinate of \(x\) is positive,
there are no edges in \(E(D,D')\), so its indicator vector vanishes.

If \(|I|\) is even, then \(D,D'\) are odd and instead
\begin{equation}
\label{geom:eq-uncross-difference}
 \chi_A+\chi_B=
 \chi_D+\chi_{D'}+2\mathbf1_{E(I,O)}.
\end{equation}
The same argument shows that \(D,D'\) are tight and that the last
indicator vanishes. Both identities hold edge by edge, including
for parallel edges. All replacement sets are nonempty and proper:
in the first case \(A\cup B\) is odd and so cannot be the even set
\(V\), and in the second case the differences are nonempty proper
subsets. In particular, if \(A\cup B=V\), only the second case can
occur. We have proved that in either case the two replacements
\(U,W\) belong to \(\mathcal S_x\) and satisfy
\begin{equation}
\label{geom:eq-uncross-tight}
 \chi_A+\chi_B=\chi_U+\chi_W.
\end{equation}

Choose an inclusionwise maximal laminar subfamily \(\mathcal L\)
of \(\mathcal S_x\) with independent cut vectors, and let \(R\)
be their span. Suppose a tight cut vector lies outside \(R\).
Among the corresponding tight sets choose \(A\) crossing as few
members of \(\mathcal L\) as possible. It crosses some
\(B\in\mathcal L\), for otherwise it could be added to the family
while preserving both laminarity and independence. Uncross \(A,B\)
as above. At least one replacement vector lies outside \(R\), since
otherwise Equation~\eqref{geom:eq-uncross-tight} and \(\chi_B\in R\)
would give \(\chi_A\in R\).

Each replacement is laminar with \(B\). Also, neither replacement
can cross a member \(C\in\mathcal L\) that did not cross \(A\).
Indeed, such a set \(C\) is laminar with both \(A\) and \(B\).
If \(C\) contains \(A\), then it meets \(B\), and laminarity
forces it to contain \(B\) as well: the other nesting direction
would imply \(A\subseteq B\). Thus \(C\) contains \(A\cup B\).
If \(C\subseteq A\), laminarity with \(B\) puts \(C\) inside
\(I\) or \(D\); it cannot contain \(B\), since \(B\nsubseteq A\).
Finally, if \(C\) is disjoint from \(A\), its relation to \(B\)
puts it inside \(D'\) or \(O\). These cases exhaust the
possibilities. Every set among \(I,A\cup B,D,D'\) is laminar with
each set in this classification.

Consequently the replacement whose vector is outside \(R\) crosses
only previously crossed members of \(\mathcal L\), and it no
longer crosses \(B\). This contradicts the minimal crossing count
of \(A\). Hence \(\mathcal L\) spans every tight cut vector.
\end{proof}

\begin{proof}[Laminar proof of Proposition~\ref{geom:prop-rank}]
Take the laminar spanning family from
Lemma~\ref{geom:lem-laminar}. Adjoin all singleton sets, if missing,
and then adjoin \(V\) as a root. The resulting inclusion tree has
\(2m\) singleton leaves. Each internal node is the disjoint union
of its children, because all singletons are present. Every proper
node has odd cardinality. A nonroot internal node therefore has an
odd number of children, at least three. The root has an even number
\(c\) of children, at least two.

If \(q\) is the number of nonroot internal nodes, the tree identity
for the number of leaves gives
\[
 2m-1
 =\sum_{S\text{ internal}}
          \bigl(\#\text{children of }S-1\bigr)
 \ge(c-1)+2q.
\]
Hence the number of proper nodes is
\(2m+q\le3m-c/2\). If \(c\ge4\), their cut vectors have rank at
most \(3m-2\) by this count. If \(c=2\), the two root children
are complementary sets, and their cut vectors coincide. The rank
is then at most the number of proper nodes minus one, again at most
\(3m-2\). The original laminar spanning family is contained among
these proper nodes. Equation~\eqref{geom:eq-rank} now proves the
claim.
\end{proof}

\subsection{The quotient-space contraction argument}

The main proof assembles compatible directions in the two contracted
graphs. The same dimension count can instead be expressed as a map to
the quotient spaces of child constraints. We include both the sharp
estimate and the weaker forms used in the entropy comparisons below.

\begin{proposition}[Supported face codimension]
\label{a:prop-supported-codimension}
For a loopless labelled multigraph on \(n=2m\ge2\) vertices and
\(x\in P(G)\), put \(A=\{e:x_e>0\}\) and \(s_A=|A|\).
Then \(s_A-d(x)\le3m-2\). In particular,
\begin{equation}
\label{a:eq-full-support-codimension}
 s_A-d(x)\le 2(n-2)\quad(n\ge4),
 \qquad s_A-d(x)=1\quad(n=2).
\end{equation}
Hence \(s_A-d(x)\le2n=4m\). For the empty graph,
\(P(G)\) is a singleton in zero coordinates and the codimension is zero.
\end{proposition}

\begin{proof}
Identify \(T(x)\) with a subspace \(U_x\) of \(\mathbb R^A\).
We prove the sharp estimate by induction on the even order \(n\).
At order two, the supported matching polytope is a simplex, whose
positive point has codimension one. For \(n\ge4\), if no nontrivial
odd cut is tight, only the degree equations constrain \(U_x\).
Thus its codimension is at most \(n\le3n/2-2\).

Otherwise choose a tight nontrivial cut \(\delta(S)\). Contract
each side in turn as in the main geometric proof, obtaining
\(x_i\in P(G_i)\), supported edge sets
\(A_i=\{e:(x_i)_e>0\}\), and direction spaces
\(U_i=T(x_i)\subseteq\mathbb R^{A_i}\). The restrictions are
feasible by the odd-cut check in the proof of
Theorem~\ref{geom:thm-polytope}. The orders satisfy
\(4\le n_i<n\) and \(n_1+n_2=n+2\).
Let \(C=A\cap\delta(S)\); this set is nonempty since \(x(C)=1\).
Every label in \(C\) occurs in both \(A_i\).

Let \(r_i:\mathbb R^A\to\mathbb R^{A_i}\) restrict coordinates,
and define
\begin{equation}\label{a:eq-quotient-map}
 \mathcal Q:\mathbb R^A\longrightarrow
  (\mathbb R^{A_1}/U_1)\oplus(\mathbb R^{A_2}/U_2),\qquad
 \mathcal Q(u)=(r_1(u)+U_1,r_2(u)+U_2).
\end{equation}
For \(u\in\ker\mathcal Q\), both restrictions are child face
directions. Choose a common positive radius such that
\(x_i+t r_i(u)\in P(G_i)\) in both signs of \(t\).
The perturbed means agree on every cut coordinate.
Lemma~\ref{geom:lem-gluing} glues them to give
\(x+tu\in P(G)\), so \(\ker\mathcal Q\subseteq U_x\).
In particular, all tight parent cuts are preserved, including those
crossing \(S\): their nonnegative slacks cannot change in opposite
signs under these two feasible perturbations.

There is one dependence between the two quotient requirements. The
cut-sum functionals
\[
 \ell_i:\mathbb R^{A_i}/U_i\longrightarrow\mathbb R,
 \qquad \ell_i(z+U_i)=\sum_{e\in C}z_e
\]
are well-defined because every child direction preserves the degree of
the contracted vertex. Each is nonzero, as the class of a unit vector
on any cut label has value one. For every \(u\in\mathbb R^A\),
\[
 \ell_1(r_1(u)+U_1)=\sum_{e\in C}u_e
                 =\ell_2(r_2(u)+U_2).
\]
Thus the image of \(\mathcal Q\) lies in the kernel of the nonzero
functional \((v_1,v_2)\mapsto\ell_1(v_1)-\ell_2(v_2)\).
This kernel has codimension one in the full quotient target.
Rank--nullity and induction now give
\begin{align*}
 s_A-d(x)&=\operatorname{codim}U_x
 \le\operatorname{codim}\ker\mathcal Q=\operatorname{rank}\mathcal Q\\
 &\le\operatorname{codim}U_1+\operatorname{codim}U_2-1\\
 &\le(3n_1/2-2)+(3n_2/2-2)-1=3n/2-2.
\end{align*}
For \(n\ge4\), this is at most \(2(n-2)\). The order-two
simplex and the empty graph give the remaining assertions.
\end{proof}

\subsection{The exact fractional residual cases}

For a general point \(x\in P(G)\), put
\[
 E_f(x)=\{e\in E:0<x_e<1\}.
\]
Let \(V_f(x)\) be the set of vertices incident to these fractional
edges, and put \(s_f=|V_f(x)|\).
The following consequence isolates the coordinates needed for the
entropy estimate.

\begin{corollary}[Fractional coordinates]
\label{a:cor-fractional-dimension}
For every \(x\in P(G)\), the integer \(s_f\) is even, and
\begin{equation}
\label{a:eq-fractional-dimension}
 d(x)\ge
 \begin{cases}
 |E_f(x)|-2(s_f-2),&s_f\ge4,\\
 |E_f(x)|-1,&s_f=2,\\
 0,&s_f=0,
 \end{cases}
 \qquad
       \sum_{e\in E_f(x)}x_e=\frac{s_f}{2}.
\end{equation}
\end{corollary}

\begin{proof}
An edge with value one exhausts the degree equation at both its
endpoints, so neither endpoint meets a fractional edge. Conversely,
at a vertex outside \(V_f(x)\), all incident coordinates are zero
or one and sum to one. Exactly one of those edges has value one.
Thus the vertices outside \(V_f(x)\) are covered by disjoint
fixed edges, and \(s_f\) is even.

Let \(H\) be the graph with vertex set \(V_f(x)\) and edge set
\(E_f(x)\), and let \(\bar x\) be the restriction of \(x\) to
\(E_f(x)\). A matching law with mean \(x\) uses every value-one edge
and no value-zero edge with probability one. Deleting the fixed
edges therefore gives a perfect matching law on \(H\) with mean
\(\bar x\), whose coordinates are all positive.
Conversely, adjoining the fixed edges to any perfect matching of
\(H\) gives a perfect matching of \(G\). This affine extension of
\(P(H)\) into \(P(G)\) sends \(\bar x\) to \(x\), and its linear part
injects the two-sided direction space at \(\bar x\) into \(T(x)\).

If \(s_f\ge 4\), Proposition~\ref{a:prop-supported-codimension} now gives
\[
 d(x)\ge \dim T(\bar x)
       \ge |E_f(x)|-2(s_f-2).
\]
If \(s_f=2\), the residual graph consists of
\(k=|E_f(x)|\ge 2\) distinguished parallel edges. Its polytope is
the simplex \(\{z\ge 0:\sum_e z_e=1\}\), and the positive point
\(\bar x\) has direction dimension \(k-1\). Hence
\(d(x)\ge k-1=|E_f(x)|-1\).
If \(s_f=0\), the point is a matching vector and the dimension
inequality is immediate; the minimal face is the singleton consisting
of that matching vector.

Finally, each vertex of \(V_f(x)\) has total fractional incident
value one. Summing these degree equations counts each fractional
edge twice, giving the second identity in
Equation~\eqref{a:eq-fractional-dimension}.
\end{proof}

\subsection{The coefficient-three and coefficient-four covariance proofs}

The coarser dimension estimates above already imply the following
entropy bounds. We use the continuity and face-law results from
Lemmas~\ref{ent:continuity} and~\ref{ent:exponential}, and the entropy
comparison in Equation~\eqref{ent:comparison}. The calculation keeps
fixed coordinates in the ambient space and takes derivatives only on
the fractional coordinates. It therefore displays directly how the
exact residual cases distinguish coefficients three and four.

\begin{samepage}
\begin{theorem}[Separable entropy bounds by covariance]
\label{a:thm-separable}
Let \(G\) be a finite loopless labelled multigraph of even order
\(n=2m\), with a nonempty perfect-matching family. With natural
logarithms and continuous endpoint conventions, define
\[
 J_c(x)=F(x)-cB(x)
   =\sum_{e\in E}\bigl[-x_e\ln x_e+c(1-x_e)\ln(1-x_e)\bigr].
\]
For each \(c\in\{3,4\}\) and every \(x\in P(G)\),
\begin{equation}\label{a:eq-separable-bounds}
 J_c(x)\le H(x)\le F(x)+B(x).
\end{equation}
For \(n=0\), all displayed quantities are zero.
\end{theorem}
\end{samepage}

\begin{proof}
The empty graph gives the stated trivial case. Assume \(n\ge2\).
For the upper bound, regard a matching law with mean \(x\) as a law
on \(\{0,1\}^E\), giving zero probability to nonmatching vectors.
Compare it with independent Bernoulli coordinates of respective
means \(x_e\).  This product law is positive wherever the matching
law is positive: a coordinate with mean zero or one is necessarily
fixed at that value under both laws.
Equation~\eqref{ent:comparison} now gives the upper bound.

For the lower bound, fix either \(c=3\) or \(c=4\).
Lemma~\ref{ent:continuity} and
continuity of \(J_c\) imply that \( H-J_c\) attains its minimum
on the compact set \(P\).
At an integral mean the only possible law is the point mass at that
matching vector, so \( H=J_c=0\).
It suffices to rule out a fractional minimizer.

Suppose, then, that \(x\) is such a minimizer.
Take its positive maximizing law \(p\) and a parameter \(\theta\)
from Lemma~\ref{ent:exponential}.
For any \(v\in\mathbb R^E\) and \(t\in\mathbb R\), define
\[
  p_a(t)=
  \frac{\exp((\theta+tv)\cdot a)}
       {\sum_{b\in\mathcal A_x}\exp((\theta+tv)\cdot b)},
  \qquad
  x(t)=\sum_{a\in\mathcal A_x}p_a(t)a.
\]
The same lemma shows that \(x(t)\) stays in the relative interior
of \(F_x\) and
\begin{equation}
   H(x(t))
  =\Phi(t)-(\theta+tv)\cdot x(t),
  \qquad
  \Phi(t)=\ln\sum_{a\in\mathcal A_x}
                       \exp((\theta+tv)\cdot a).
  \label{a:eq-tilted-entropy}
\end{equation}
Thus we may differentiate along these curves without assuming
differentiability of \( H\) across faces.

Let
\[
  C=\sum_{a\in\mathcal A_x}p_a(a-x)(a-x)^{\mathsf T}
\]
be the covariance matrix at \(t=0\).
Differentiating the finite sums gives
\(p_a'(0)=p_a\,v\cdot(a-x)\), hence \(x'(0)=Cv\).
For any \(w\in\mathbb R^E\),
\[
  w^{\mathsf T}Cw
  =\sum_{a\in\mathcal A_x}p_a\bigl(w\cdot(a-x)\bigr)^2.
\]
All \(p_a\) are positive, and
\(\operatorname{span}\{a-x:a\in\mathcal A_x\}=T(x)\).
Since \(C\) is symmetric, it follows that
\begin{equation}
  \ker C=T(x)^\perp,\qquad
  \operatorname{im}C=T(x),\qquad
  \operatorname{rank}C=d(x).
  \label{a:eq-covariance-face-space}
\end{equation}
At a fractional mean, some coordinate has variance
\(C_{ee}=x_e(1-x_e)>0\), so \(d(x)\ge1\).
Also \(x''(0)\in T(x)\), because the whole curve lies in
\(\operatorname{aff}F_x=x+T(x)\).

On the fractional coordinates, write
\[
  g_e=-\ln x_e-c\ln(1-x_e)-(c+1),
  \qquad
  D_{ee}=-\frac1{x_e}+\frac c{1-x_e};
\]
let \(D\) be diagonal, and set \(g_e=D_{ee}=0\) on the other
coordinates.  These entries represent the derivatives of \(J_c\)
restricted to the face.
Indeed, a coordinate with \(x_e\in\{0,1\}\) has that value on every
face atom, so its contribution to \(J_c(x(t))\) is constantly zero.
A fractional coordinate takes both bit values among the face
atoms; positivity of \(p(t)\) keeps \(0<x_e(t)<1\) for every finite
\(t\).  Consequently the usual chain rule on fractional coordinates
gives all derivatives of \(J_c(x(t))\).

From \(\Phi'(t)=v\cdot x(t)\) and
Equation~\eqref{a:eq-tilted-entropy}, we obtain
\[
  \frac{d}{dt} H(x(t))
  =-(\theta+tv)\cdot x'(t).
\]
Both signs of \(t\) are permitted.  The minimum property at \(t=0\)
therefore gives
\[
  (\theta+g)\cdot Cv=0
  \qquad\text{for every }v\in\mathbb R^E.
\]
By Equation~\eqref{a:eq-covariance-face-space}, \(\theta+g\) is
orthogonal to \(T(x)\), and in particular
\((\theta+g)\cdot x''(0)=0\).
The second derivative at the minimum is accordingly
\begin{equation}
\begin{aligned}
  0
  &\le
    \left.\frac{d^2}{dt^2}
      \bigl( H(x(t))-J_c(x(t))\bigr)\right|_{t=0}\\
  &=-v^{\mathsf T}Cv-(\theta+g)\cdot x''(0)
                     -(Cv)^{\mathsf T}D(Cv)\\
  &=-v^{\mathsf T}Cv-(Cv)^{\mathsf T}D(Cv).
\end{aligned}
\label{a:eq-entropy-second-variation}
\end{equation}

Let \(u_1,\ldots,u_d\) be orthonormal eigenvectors for the positive
eigenvalues \(\lambda_1,\ldots,\lambda_d\) of \(C\), where
\(d=d(x)\).
Substituting \(v=u_i\) in
Equation~\eqref{a:eq-entropy-second-variation} and dividing by
\(\lambda_i>0\) gives
\(\lambda_i u_i^{\mathsf T}(-D)u_i\ge1\).
Summing and using
\(C=\sum_{i=1}^d\lambda_i u_i u_i^{\mathsf T}\), we find
\begin{equation}
  d(x)\le
  \sum_{i=1}^d\lambda_i u_i^{\mathsf T}(-D)u_i
  =\operatorname{tr}(C(-D)).
  \label{a:eq-entropy-trace-lower}
\end{equation}
Only the spectral decomposition of \(C\) is used here; no inverse
of \(C\), or simultaneous diagonalization of \(C\) and \(D\),
is needed.

Each matching coordinate is Boolean, so
\(C_{ee}=x_e(1-x_e)\).
Writing \(s_f\) for the number of vertices incident with
\(E_f(x)\), Corollary~\ref{a:cor-fractional-dimension} gives
\begin{equation}\label{a:eq-covariance-trace-c}
\begin{aligned}
 \operatorname{tr}(C(-D))
 &=\sum_{e\in E_f(x)}x_e(1-x_e)
       \left(\frac1{x_e}-\frac c{1-x_e}\right)\\
 &=|E_f(x)|-(c+1)\sum_{e\in E_f(x)}x_e
   =|E_f(x)|-\frac{c+1}{2}s_f.
\end{aligned}
\end{equation}
For \(c=4\), the weaker consequence
\(d(x)\ge|E_f(x)|-2s_f\) of the residual dimension bound suffices:
\[
 \operatorname{tr}(C(-D))=|E_f(x)|-\tfrac52s_f
 <|E_f(x)|-2s_f\le d(x),
\]
because a fractional point has \(s_f>0\).
For \(c=3\), the trace is \(|E_f(x)|-2s_f\), so we use the exact
residual cases. If \(s_f\ge4\), then
\[
 d(x)\ge|E_f(x)|-2(s_f-2)
       =|E_f(x)|-2s_f+4
       >\operatorname{tr}(C(-D)).
\]
If \(s_f=2\), then
\[
 d(x)\ge|E_f(x)|-1>|E_f(x)|-4
       =\operatorname{tr}(C(-D)).
\]
These exhaust the positive even values of \(s_f\). For either
coefficient, the result contradicts
Equation~\eqref{a:eq-entropy-trace-lower}. Thus no fractional minimizer
exists; the minimum is attained at an integral mean and equals zero.
This proves the lower bound for each coefficient by its own covariance
contradiction.
\end{proof}

\subsection{Two other entropy comparisons}

The covariance comparisons just proved give separable counting estimates
and a short scalar proof of the nonlinear bound. We give those deductions
here before the distinct normalized-covariance proof in
Appendix~\ref{b:section}. As in Section~\ref{cons:section},
$N(G)$ denotes the number of labelled perfect matchings, and
$B(x)\le m$ by Equation~\eqref{cons:B-bound}.

\begin{corollary}[Separable and nonlinear consequences]\label{cons:legacy-defect}
For every $x\in P(G)$ with $|V(G)|=n=2m>0$,
\begin{align}
 F(x)-4B(x)&\le F(x)-3B(x)\le H(x)\le F(x)+B(x),\label{cons:legacy-separable}\\
 F(x)-H(x)&\le\min\{F(x),3m\}
                  \le8m\bigl(1-e^{-F(x)/m}\bigr).\label{cons:legacy-nonlinear}
\end{align}
Furthermore,
\begin{equation}\label{cons:legacy-bin-gap}
 \ln N(G)\le\max_{x\in P(G)}[F(x)+B(x)]
       \le\ln N(G)+2n
       \le\ln N(G)+\tfrac52n.
\end{equation}
\end{corollary}
\begin{proof}
Theorem~\ref{a:thm-separable} proves
$F-3B\le H\le F+B$ by the contraction and covariance route; its
coefficient-four proof gives $F-4B\le H$ as well. Since $B\ge0$, the
first two terms are ordered as displayed. Nonnegativity of entropy and
$B\le m$ give $F-H\le\min\{F,3m\}$.
Put $t=F(x)/m\ge0$ and $f(t)=8(1-e^{-t})$. On $[0,\ln2]$, one has
$f(0)=0$ and $f'(t)=8e^{-t}\ge4\ge1$, hence $f(t)\ge t$.
For $t\ge\ln2$, one has $f(t)\ge4\ge3$.
Thus $\min\{t,3\}\le f(t)$, proving the nonlinear inequality including
$t=0$. This deduction is independent of the normalized inverse-covariance
proof of Theorem~\ref{b:thm-nonlinear}.

For the optimized bound, $H\le\ln N$ and
$(F+B)-(F-3B)=4B\le4m=2n$ give the upper error $2n$.
Using $F-4B$ instead gives the separate original comparison
$(F+B)-(F-4B)=5B\le5m=5n/2$. The uniform matching law supplies the
lower bound exactly as in Corollary~\ref{cons:entropy-count}.
The sharper error $3m-2$ in that corollary improves both comparisons
on the same domain.
\end{proof}

\subsection{Structural estimates from the other entropy bounds}

For $k\ge1$, put $b_1=1$ and
$b_k=k(1-1/k)^{3(k-1)}$ for $k\ge2$.

\begin{corollary}[Coefficient-three and nonlinear counts]\label{cons:legacy-counts}
A loopless labelled multigraph $G$ on $2m>0$ vertices containing $k\ge1$
pairwise edge-disjoint perfect matchings satisfies
\begin{align}
 N(G)&\ge b_k^m\ge(e^{-3}k)^m,\label{cons:legacy-three-count}\\
 N(G)&\ge\bigl[k\exp\{-8(1-1/k)\}\bigr]^m
                         \ge(e^{-8}k)^m.\label{cons:legacy-eight-count}
\end{align}
The same two conclusions hold when $G$ is $k$-regular and every odd cut
has at least $k$ edges.\label{cons:legacy-regular}
\end{corollary}
\begin{proof}
For the selected matching mixture in Theorem~\ref{cons:structural},
$F(x)=m\ln k$ and $x$ vanishes outside their labelled union.
For $k\ge2$, $F(x)-3B(x)=m\ln k+3m(k-1)\ln(1-1/k)$.
Using $H(x)\le\ln N(G)$ and the coefficient-three inequality in
Corollary~\ref{cons:legacy-defect}, then exponentiating, gives the first
bound. The inequality $(k-1)\ln(1-1/k)\ge-1$ from the proof of
Theorem~\ref{cons:structural} gives its coarse form. At $k=1$ the selected matching gives $N(G)\ge1=b_1^m$.

At this same mean, the nonlinear inequality gives
$\ln N(G)\ge m\ln k-8m(1-1/k)$, including $k=1$.
Exponentiating and using $1-1/k\le1$ proves the second line.
For a regular graph satisfying the odd-cut hypothesis,
Corollary~\ref{cons:regular} proves that the uniform vector is feasible;
the same evaluations then prove both statements without requiring
a partition of the edges into perfect matchings. Its cut-parity argument also proves the
$(k-1)$-edge-connected even-order case.
\end{proof}

For $k\ge2$, the finite-order bound of Theorem~\ref{cons:structural}
implies the coefficient-three bound directly, since $0<1-1/k<1$ and
$2(m-1)(k-1)\le3m(k-1)$. The coefficient-three logarithmic loss is at
most $3$, while $8(1-1/k)\ge4$ for $k\ge2$, so it also dominates the
second displayed base. We have nevertheless derived that base from
the nonlinear inequality to retain the distinct implication. At $k=1$,
the refined bounds $b_1^m$ and $[k\exp\{-8(1-1/k)\}]^m$ both equal one.

\begin{corollary}[An exact partition of the ambient edges]\label{cons:legacy-partition}
If the entire edge set of a loopless labelled graph on $2m>0$ vertices
is partitioned by $k\ge1$ perfect matchings, then
\[
 N(G)\ge\bigl[k\exp\{-8(1-1/k)\}\bigr]^m\ge(e^{-8}k)^m.
\]
When $k=1$ the count is exactly one. If $k=0$ and the order is positive,
the graph is edgeless and its count is zero; only the coarse bound is
then asserted. On the empty vertex set there is one empty matching.
\end{corollary}
\begin{proof}
The partition provides the hypothesis of Corollary~\ref{cons:legacy-counts};
its uniform law has marginal $1/k$ on every ambient edge. Thus the
nonlinear proof just given applies verbatim. For $k=1$ the entire graph
is a single perfect matching. For $k=0$ it has no edges. The empty case
follows directly from the definition of a matching and does not use
expressions containing $1/m$.
\end{proof}

The exact-partition hypothesis in the last corollary is stronger than
containing a selected union. In particular, a graph containing one
perfect matching may have more than one, whereas the exact one-matching
union has count one.
